\documentclass{article}

%% Packages %%
\usepackage[utf8]{inputenc}
\usepackage{graphicx}
\usepackage{hyperref}
\usepackage[dvipsnames]{xcolor}

%% Commands %%

% Define new commands to store the trim values
\newcommand{\trimLR}{2cm}
\newcommand{\trimtop}{3.5cm}
\newcommand{\trimbottom}{4cm}

% This command creates a figure environment with a specific image:
% \myfigure{page}{trim_width}{trim_height}{filename}
%
% Arguments:
%  page: The page number of the image in the file.
%  trim_width: The width of the trim box.
%  trim_height: The height of the trim box.
%  filename: The name of the file that contains the image.
%
% Usage:
%  \myfigure{1}{1cm}{1cm}{image.png}
% This will create a figure with the image from page 1 of image.png, trimmed to 1cm by 1cm.
\newcommand{\myfigure}[4]{
  \begin{figure}[h]
    \centering
    \includegraphics[page=#1, trim = #2 #3 #2 \trimtop, clip, width=\textwidth]{#4}
  \end{figure}
}

% Sets the length of the vertical space between paragraphs to 8pt.
\setlength{\parskip}{8pt}

%% Document %%

\title{Business Information Systems - Part 2}
\author{Gabriele Lazzarelli}

\begin{document}

% Debugging purposes
\pagecolor{Goldenrod}

\maketitle
\tableofcontents

% \section{Exam Schedules and
    Evaluations}\label{exam-schedules-and-evaluations}

Students in group two should attend the lecture, while students in group
one should study the subject and the lecture slides for the exam. Group
two students will be evaluated in class, while group one students will
have an exam based on the material covered in the lecture. Some students
from group one have asked about the midterm exam for BIS two, and the
answer is yes, there will be a midterm exam. The date for the final exam
can be found on the course schedule on the WEB platform, and it will
only cover the material from BIS two. The midterm exam for BIS one has
already taken place, and there will be a separate midterm exam for BIS
two.
During official exam dates, excluding midterms, both BIS one and BIS two
exams are usually scheduled at different times. However, starting from
September, when there are fewer students, the exams are held
simultaneously in the same room. In June and July, BIS one is typically
conducted first, followed by BIS two after an hour. I will inform you of
the specific instructions via email when the time comes.
% \section{Workforce Management Lecture}\label{workforce-management-lecture}

\myfigure{2}{\trimLR}{4cm}{images/00 - WFM.pdf}

\subsection{Extended ERP and Workforce Management}\label{extended-erp-and-workforce-management}

Now, let's move on to the topic of workforce management. Extended ERP is
the focus here. Extended ERP aims to integrate all information flows and
processes involving not only internal actors like employees but also
external players such as clients and partners of the company. This
comprehensive approach ensures efficient interactions between the
company and its stakeholders.

\myfigure{3}{\trimLR}{3.5cm}{images/00 - WFM.pdf}

\subsection{Client Relationships and Transaction Models}\label{client-relationships-and-transaction-models}

Workforce management is closely tied to the relationship between a
company and its clients, particularly during the post-sale phase. This
relationship is governed by a transaction model, which we discussed in
the context of transaction theory. It's important to note that there are
various types of transactions, depending on the nature of the goods or
services involved. However, all transactions ideally consist of four
phases, with the final phase being post-settlement or post-sale
services. Workforce management focuses on these post-sale services and
the ongoing relationship between the company and the client.


\subsection{Post-Sale Services and WFM in Manufacturing}\label{post-sale-services-and-wfm-in-manufacturing}

When discussing workforce management, it is often in the context of the
manufacturing industry. In this industry, companies produce physical
products that are installed at the customer's location. Examples of
these installations include elevators, household appliances like washing
machines and dishwashers, as well as appliances for electricity or gas.
The focus of workforce management in this context is on maintaining
these installations.

Once a customer has purchased an appliance, they may encounter issues
that require maintenance. The manufacturer is responsible for addressing
these maintenance requests. Typically, the customer has already paid for
the appliance, so the transaction has been completed. However,
additional services are needed to resolve any issues with the product.
To provide these services, the company must send specialized employees
to the customer's location.

\subsection{Maintenance Challenges and
    Costs}\label{maintenance-challenges-and-costs}

For example, let's say your elevator is broken. In this scenario,
someone from the company needs to come to your building to fix it.
However, having a workforce that needs to travel to different locations
comes with a significant cost for the company. It would be much more
cost-effective if the client could bring the appliance to the company
for repairs.

To illustrate this, let's consider the example of a broken cell phone.
When your cell phone breaks, you don't have someone come to your place
to fix it. Instead, you go to a shop that specializes in cell phone
repairs. You might go to a shop in Chinatown, for example, where they
provide excellent service. You drop off your phone, go about your day,
and then return to pick it up once it's fixed. This process requires you
to invest your time in physically going to the shop, which can be a
significant amount of time. However, you are not charged for the
traveling costs.

In contrast, when a company needs to send a workforce to your home to
fix a broken appliance, it becomes more complex. The company has to
consider the costs associated with sending their workforce to your
location, including travel expenses and the time required for travel.
These costs are eventually factored into the price of the service, which
you, as the client, will have to pay. Unlike when you take your cell
phone to a shop, where you are not charged for traveling costs, the
company needs to include these costs in their service fees.
Additionally, the company must ensure that they send the right skilled
workers to address the specific issue you are facing.

\subsection{Emergency vs.~Routine
    Maintenance}\label{emergency-vs.-routine-maintenance}

\myfigure{4}{\trimLR}{7cm}{images/00 - WFM.pdf}

And so, when it comes to fixing appliances like elevators, the team
responsible for the repairs needs to have the necessary skills and
competencies. Unlike other appliances that can be taken to the company
for repairs, elevators need to be fixed on-site, which often requires
multiple visits. The first visit is to assess the problem, and
subsequent visits are made with the appropriate team and materials to
fix the issue. This process is complex and involves managing response
times, which is crucial for customer satisfaction and safety.

In the case of elevators, response time is particularly important
because a malfunctioning elevator can pose a safety risk, especially if
someone is trapped inside. For example, if a patient in a hospital is
being transferred between departments and the elevator gets stuck, it
becomes a matter of urgency to rescue them. The level of service
provided in these situations is not only important for customer
satisfaction but also for ensuring safety. Incidents where people are
trapped in elevators and their lives are at risk are rare but
significant. Therefore, the level of service and efficiency in handling
emergencies is crucial.

Maintenance plays a vital role in assessing the quality of service
provided by a company. When an appliance breaks down, it is an
opportunity for the company to demonstrate its commitment to customer
satisfaction. If the company efficiently handles the maintenance
process, customers are more likely to remain loyal. Studies have shown
that satisfied customers who receive efficient maintenance services are
more likely to continue purchasing products from the same company.
Acquiring new customers is challenging, so maintaining the loyalty of
existing customers is essential for long-term competitiveness.

Elevators are a prime example of the importance of maintenance because
they are long-lasting investments. Instead of replacing a broken
elevator, it is more cost-effective to repair it. Therefore, maintenance
becomes even more critical for elevators. However, maintenance requests
can be divided into two types: routine maintenance and emergency
maintenance. Routine maintenance, such as changing the oil in a car
engine, can be planned and scheduled in advance. It is preventive
maintenance that helps avoid breakdowns and can be optimized for
efficiency. On the other hand, emergency maintenance is unpredictable
and requires immediate attention.

\myfigure{5}{\trimLR}{11cm}{images/00 - WFM.pdf}

Routine maintenance is relatively easy to manage and profitable for
companies. It can be planned ahead of time, and the effort required is
known in advance. This type of maintenance is designed to be
straightforward and can be performed without specialized expertise or
spare parts. It is a predictable demand that allows companies to
allocate resources efficiently. In contrast, emergency maintenance is
unpredictable and requires a reactive response.

In summary, the composition of the workforce and their competencies are
crucial for fixing appliances like elevators. Response time is essential
for customer satisfaction and safety, especially in emergency
situations. Maintenance is an opportunity for companies to demonstrate
their commitment to customer service and build loyalty. Routine
maintenance can be planned and optimized, while emergency maintenance
requires immediate attention. By understanding the different types of
maintenance and their importance, companies can ensure long-term
competitiveness and customer satisfaction.

\subsection{Customer Service and Maintenance
    Requests}\label{customer-service-and-maintenance-requests}

In the case of sudden problems that require management, the typical
scenario involves a client calling the company to report the issue. This
asynchronous communication requires a synchronous response from the
company to handle the problem efficiently. However, this type of demand
is difficult to forecast, leading to challenges in planning and
efficiency. Emergency maintenance, especially if it involves safety,
must be prioritized regardless of location. This means that companies
must be prepared to serve clients in remote areas, even if it is less
convenient or profitable.

\myfigure{8}{\trimLR}{5cm}{images/00 - WFM.pdf}

For manufacturing companies, the transaction with the customer typically
ends with the payment, and there is no ongoing contract or relationship.
This can lead some clients, especially businesses, to seek maintenance
services from other providers instead of the manufacturer. This is often
driven by the perception that independent maintenance services are
cheaper than going to the manufacturer. Smaller maintenance services
tend to focus on profitable areas with high concentrations of potential
issues, while leaving less profitable areas to the manufacturer. As a
result, manufacturers are left with the more challenging and less
profitable emergency maintenance tasks.

\myfigure{6}{\trimLR}{6cm}{images/00 - WFM.pdf}

Emergency maintenance is generally not profitable, especially when it
requires sending a workforce to the customer's location. Maintenance
activities can be costly due to the geographical distribution of
customers, the expense of visits, and the difficulty of finding
employees who are skilled at problem-solving. Employees who excel at
fixing appliances may even leave to start their own maintenance
services, becoming competitors to the manufacturer. Additionally,
handling maintenance requests can be challenging, as customers often
express their frustration and anger without providing essential
information about the appliance and its model.

Overall, routine maintenance is typically more profitable than emergency
maintenance. However, the nature of emergency maintenance, with its
generic and often emotionally charged requests, makes it more difficult
to manage efficiently.

\subsection{Optimizing the Maintenance
    Process}\label{optimizing-the-maintenance-process}

To optimize the maintenance process, it is important for customers to provide essential information about the appliance, such as its location and details. However, many customers are not aware of this and may struggle to accurately describe the problem. This is where the call center operators play a crucial role in asking the right questions to gather the necessary information. Unfortunately, level one call center operators often lack the skills to ask precise questions, which can lead to inefficiencies in the process.

\myfigure{9}{\trimLR}{4cm}{images/00 - WFM.pdf}

To address this issue, it is important for organizations to ensure that the call center operators have access to an enterprise resource planning (ERP) system. This system can help them set up appointments and share the workforce's agenda, ensuring that the right resources are available when needed. Additionally, having spare parts readily available is essential to avoid delays in the maintenance process. However, tight service level agreements (SLAs) can pose challenges, as the workforce needs to be adequately sized to handle peak calls, which can increase costs.

\myfigure{7}{\trimLR}{7.5cm}{images/00 - WFM.pdf}

It is common for maintenance to be seen as a non-profitable and challenging process. However, it is crucial to change this mindset and view maintenance as a service that can help build customer loyalty. Outsourcing maintenance may seem like a good solution, but it can have negative repercussions on brand equity if issues arise. Therefore, organizations should carefully consider the level of service provided by maintenance service providers before outsourcing.

The worst possible maintenance process is one where the customer provides limited information about the issue, leading to multiple visits and coordination challenges within the workforce. This not only inconveniences the customer but also increases costs. Additionally, missed opportunities for cross-selling\footnotemark{} can occur during maintenance interventions. This is the perfect time to offer maintenance subscriptions or additional services to customers.

To optimize the maintenance process, it is crucial for call center operators to ask the right questions and gather accurate information from customers. This will help streamline the process, reduce costs, and improve customer satisfaction.

\footnotetext{Cross-selling is a sales strategy where a company encourages customers to purchase additional products or services related to their initial purchase. The goal of cross-selling is to increase the average transaction value and maximize revenue from each customer. This strategy involves suggesting complementary or supplementary items that go hand-in-hand with the customer's original purchase, thereby enhancing their overall experience and meeting more of their needs.}

\subsection{Knowledge Management and Business
    Intelligence}\label{knowledge-management-and-business-intelligence}

\myfigure{10}{\trimLR}{3cm}{images/00 - WFM.pdf}

To ensure efficient maintenance processes, it is important to have an
intelligent Q\&A system in place. This system allows the maintenance
workforce to diagnose problems accurately and quickly. By diagnosing the
problem, the necessary spare parts can be prepared in advance. The ERP
system can be used to check if the spare parts are available in the
physical warehouse and ensure they are loaded onto the truck. This way,
the technician can make an appointment at the right time and avoid any
delays.

Invoicing is another crucial aspect of the maintenance process. It is
essential to make the customer aware of why they are spending money and
provide a breakdown of the costs, including the price of the spare
parts. By doing this, the customer can understand the value they are
receiving for their investment. It is recommended to print the invoice
as soon as possible, preferably during the discussion with the customer,
to ensure accuracy and transparency.

Training the workforce to handle the administrative part of the job is
necessary. Although technicians may not enjoy administrative tasks, it
is important to equip them with the skills to handle invoicing and other
administrative responsibilities. This can be facilitated by using a
simple application that tracks materials used during maintenance
services and simplifies the invoicing process.

\myfigure{11}{\trimLR}{6cm}{images/00 - WFM.pdf}

Implementing knowledge management is crucial for an intelligent Q\&A
system. By analyzing past maintenance interventions and building
diagnostic intelligence, the system can ask the right questions and
provide accurate solutions. Additionally, appliances equipped with
sensors can self-diagnose and provide error codes, making the Q\&A
process even more efficient.

Furthermore, by analyzing data from maintenance services, companies can
gain valuable insights into common issues with different appliance
models. This information can be used to improve product design and solve
recurring problems. However, it is important to note that using
knowledge management to intentionally shorten the lifespan of appliances
for profit is unethical and should not be practiced.

In summary, the key actions to achieve an ideal maintenance process
include implementing knowledge management, using an intelligent Q\&A
system, training the workforce in administrative tasks, and analyzing
data for continuous improvement.

\subsection{Case Studies}\label{case-studies}

\subsubsection{OTIS Elevators Example}\label{otis-elevators-example}

Knowledge management is essential for analyzing data and understanding
maintenance issues. By collaborating with the maintenance team and Research and Development (R\&D) team, you can identify and fix problems before they become emergencies. This
not only helps avoid costly interventions but also allows for the
transformation of emergency maintenance into routine profitable
maintenance. By analyzing information on maintenance interventions and
working with R\&D, you can discover the most common and recurring issues
and implement preventive maintenance to avoid emergencies. This
knowledge can be kept confidential and used to offer a competitive
service in the market.

\myfigure{13}{\trimLR}{7cm}{images/00 - WFM.pdf}
\myfigure{14}{\trimLR}{8cm}{images/00 - WFM.pdf}

A case study that exemplifies these issues is the OTIS Elevator company. They
faced the challenge of unprofitable maintenance services due to
competition. Competitors focused on the easiest and most profitable
interventions in urban areas, leaving OTIS Elevator with elevators in
remote areas that lacked maintenance subscriptions and routine
maintenance. This resulted in numerous emergency calls. Additionally,
OTIS Elevator mismanaged interventions and faced payment issues from
clients. While a small percentage of unpaid interventions can be
managed, if it grows, it becomes a problem as it affects the
profitability of the maintenance processes.

\myfigure{15}{1.5cm}{5cm}{images/00 - WFM.pdf}

To address these challenges, OTIS Elevator decided to change their
maintenance processes and implement knowledge management. This involved
conducting proper business intelligence analysis and data analytics to
uncover the root causes of issues and their relationships. By working
with data scientists and R\&D, they established routine maintenance
interventions that reduced the need for emergency interventions. They
also leveraged their ability to diagnose issues before they occur and
offered highly competitive prices for their maintenance services.

By understanding the issues and their causes, OTIS Elevator was able
to proactively replace parts after three years, resulting in cost
savings for customers and the ability to offer maintenance interventions
at a discounted price compared to the market. Through this process, they
invested in customer relationships and subscription maintenance
services, reducing the need for emergency maintenance by 30 percent and
becoming more competitive in terms of pricing. This demonstrates the
benefits of industrializing processes by standardizing and making them
more efficient.

By utilizing knowledge management and business intelligence, OTIS
Elevator was able to act proactively and avoid problems. This approach
not only improved their maintenance processes but also enhanced their
overall service quality.

\myfigure{16}{\trimLR}{4cm}{images/00 - WFM.pdf}
And then they implemented an industrialized process. The value
proposition became subscribing to their maintenance service, which would
result in a 30 percent reduction in emergency maintenance interventions
and cost savings for customers. This would also improve the quality of
their lives by minimizing the need to handle emergencies. By offering
this service, they could shift the blame to competitors if any issues
arose due to improper routine maintenance. In business contexts, people
are willing to pay to avoid problems, making this a strong value
proposition.

\myfigure{17}{\trimLR}{5cm}{images/00 - WFM.pdf}

They successfully sold many maintenance services by
investing in IoT and offering free sensors and diagnostic tools for old
elevators. This allowed them to send the right technicians with the
necessary spare parts at the right time. Ultimately, their solution
extended the lifespan of the equipment and protected the owner's
investment. They positioned maintenance as a service rather than a
costly and unprofitable process.

\myfigure{18}{\trimLR}{6cm}{images/00 - WFM.pdf}

Their web-based system included remote
monitoring capabilities, allowing them to proactively address issues
before they occurred. They would contact customers to schedule
maintenance visits, creating a high level of trust and customer
satisfaction. This approach allowed them to charge premium prices,
similar to a luxury hotel experience.

\myfigure{19}{\trimLR}{11cm}{images/00 - WFM.pdf}

\subsubsection{Utility Company and Smart
    Meters}\label{utility-company-and-smart-meters}

\myfigure{21}{\trimLR}{9cm}{images/00 - WFM.pdf}

Let's discuss the use of intelligent meters by utility companies. These
meters provide real-time information on consumption, allowing for
accurate billing without the need for manual readings. Instead of
sending a workforce to check the meters, utility companies can position
intelligent meters that provide consumption data in real time. This
eliminates the need for manual readings and reduces costs.

\myfigure{26}{\trimLR}{7cm}{images/00 - WFM.pdf}

However, utility companies have not fully utilized the potential of the
information provided by these meters. Currently, they mainly use the
data for self-diagnosis and fixing any issues with the meters. While
this is useful, it does not offer a compelling value proposition to
customers. The meters themselves are simple and do not require frequent
maintenance like elevators, so the potential for additional services is
limited.

The value proposition of these meters is primarily focused on improving
invoicing accuracy. However, utility companies are now reconsidering
their approach due to sustainability concerns and the increasing cost of
energy. They are exploring the idea of using renewable energy sources
and optimizing energy consumption to reduce costs. This could
potentially create a more compelling value proposition for customers.

In conclusion, while utility companies have not fully capitalized on the
potential of intelligent meters, there is ongoing exploration of new
value propositions. The example of the OTIS elevator remains a
favorite case study, highlighting the importance of integrating
maintenance into the overall service cycle.


% \section{KPMG on ERP Systems}\label{technical-presentation}

\subsection{Introduction and Agenda}\label{introduction-and-agenda}

\begin{figure}[!h]
    \centering
    \includegraphics[page=2, trim = 2cm 7cm 2cm 0cm, clip, width=\textwidth]{images/01 - KPMG_IT Advisory Services_ENG_v1.pdf}
\end{figure}

We have plenty of time this morning, so let's make the most of it. I
propose structuring the next 60 minutes as follows:

First, I will provide a brief introduction about myself and the company
I work for. Then, we will dive into the documentation I have prepared.

To begin, my name is Marco Trammelli, and I am a senior manager of advisory at the KPMG company. I will give you an overview of the system we
will be discussing today, whic  h serves as a prerequisite for
understanding the rest of the session.

Next, we will move on to the second part of today's session, where I
will introduce a test case scenario that you will be working on in the
coming days. After that, I will hand over the stage to my colleague from
the HR department.

\subsection{System Overview}\label{system-overview}

\subsubsection{Historical Perspective and
    Evolution}\label{historical-perspective-and-evolution}

Let's start with a brief introduction to the system, considering its
historical perspective and evolution. We'll also provide an overview of
a sample system, the SAP ERP system, which is widely used in the
market. If time permits, we'll then move on to our recommendations for
implementation and development projects.

\begin{figure}[!h]
    \centering
    \includegraphics[page=3, trim = 2cm 5cm 1cm 0cm, clip, width=\textwidth]{images/01 - KPMG_IT Advisory Services_ENG_v1.pdf}
\end{figure}

The term ``legacy system'' refers to a collection of autonomous
applications, each with its own database. These systems are not
interconnected, and interfaces need to be developed and implemented to
establish communication between them. In the following roadmap, we'll
highlight the evolution of these systems over the decades, starting from
the first enterprise resource planning system in the 1970s.

\subsubsection{Legacy Systems vs.~ERP
    Systems}\label{legacy-systems-vs.-erp-systems}

\begin{figure}[!h]
    \centering
    \includegraphics[page=4, trim = 2cm 3cm 1cm 0cm, clip, width=\textwidth]{images/01 - KPMG_IT Advisory Services_ENG_v1.pdf}
\end{figure}

Basically, this system was initially developed for manufacturing and
production processes within the organization. Over time, additional
features were added to support other functions and processes. This
evolution led to what we now call the extended ERP system.

\begin{figure}[!h]
    \centering
    \includegraphics[page=5, trim = 2cm 3cm 2cm 0cm, clip, width=\textwidth]{images/01 - KPMG_IT Advisory Services_ENG_v1.pdf}
\end{figure}

Now, let's briefly compare legacy systems with ERP systems. Legacy
systems are often referred to as closed systems. They store information
about past actions in their own databases. Each function in the
organization, such as purchasing, shipping, production, warehouse
administration, and sales, typically has its own separate system. This
lack of a common repository can lead to duplicated activities and
operational inconsistencies across different business functions.

\begin{figure}[!h]
    \centering
    \includegraphics[page=9, trim = 2cm 2.5cm 2cm 0cm, clip, width=\textwidth]{images/01 - KPMG_IT Advisory Services_ENG_v1.pdf}
\end{figure}

On the other hand, the ERP system is a
unique system that can simulate and propose future scenarios instead of
focusing on the past. This system shares a common database, which I will
refer to as a table or database. Additionally, there is a sharing of
operational logic among different functional areas within a company,
such as purchasing, shipping, and production. Each block in this diagram
represents a specific function within the organization. It is important
to note the close interconnection of activities carried out in these
functional areas. Instead of developing separate programs for each
function, the system allows for the development of shared programs for
multiple functions. These key features cover various business processes,
including logistics, accounting, production, and human resources. These
processes are supported by a set of application elements known as
modules. Throughout this presentation, we will explore how each module
represents a specific function within the organization.

\begin{figure}[!h]
    \centering
    \includegraphics[page=18, trim = 2cm 4cm 2cm 0cm, clip, width=\textwidth]{images/01 - KPMG_IT Advisory Services_ENG_v1.pdf}
\end{figure}

Can someone explain what a legacy system is and what the main
differences are between a legacy system and an ERP system? Let's imagine a
company and focus on the purchase processes, such as purchase
requisitions, purchase orders, receiving invoices, etc. Invoicing is
related to accounting and finance, while purchasing is mainly related to
the buyer or someone in the purchasing department. It's not interesting
to discuss the connection between finance and accounting in these
processes. A legacy system uses different tables that need to be
reconciled, requiring the implementation of multiple interfaces and
effort to check for redundant data. On the other hand, an ERP system has
a common structure with a unique table, eliminating the need for
reconciliation and interfaces. By implementing an ERP system from the
beginning, an organization can focus on developing efficiency in
processes instead of spending time and money on reconciling data. Legacy
systems are difficult to integrate with external tools, while ERP
systems already have connections in place. ERP systems also come with
embedded best practices for accounting, finance, and purchasing
processes, making them ready to use with just configuration based on the
organization's specificity. Legacy systems require specific interfaces
to integrate with other software and require the organization to align
with best practices. If a specific process is not in line with best
practices, it cannot be included in an ERP system without revising and
adapting it. When selecting an ERP system, industry specificity,
specific organizational issues, business strategy, industrial
development plans, processes, and resources should be considered. These
factors will guide the software selection process, whether it involves
implementing an ERP system from the market or looking for specific
software applications from local vendors.

\subsubsection{ERP Implementation and
    Development}\label{erp-implementation-and-development}

\begin{figure}[!h]
    \centering
    \includegraphics[page=8, trim = 2cm 5cm 2cm 0cm, clip, width=\textwidth]{images/01 - KPMG_IT Advisory Services_ENG_v1.pdf}
\end{figure}

The logistics department, accounting, production, sales, and human
resources all benefit from the adoption of an ERP system. By sharing the
same repository and database, businesses can ensure the uniqueness and
universality of their data and business information. Major software
vendors like Microsoft and SAP produce and release ERP systems, such as
MRP (Material Requirements Planning).

From a technical perspective, ERP systems use a relational database and
operate on a client-server architecture. They also offer vertical
solutions tailored to specific industries, such as automotive or
healthcare, which have unique business processes and requirements.

Implementing an ERP system allows companies to leverage best practices
and standardized processes across departments, eliminating the need to
reinvent the wheel for each company. This saves time and resources while
ensuring efficiency and consistency.

Moving forward, many companies are already using ERP systems like
Microsoft or SAP, but they continue to seek improvements and
enhancements to their systems.

\subsubsection{Future Directions and Cloud
    Computing}\label{future-directions-and-cloud-computing}

\begin{figure}[!h]
    \centering
    \includegraphics[page=12, trim = 2cm 3cm 2cm 0cm, clip, width=\textwidth]{images/01 - KPMG_IT Advisory Services_ENG_v1.pdf}
\end{figure}

In recent years, there have been common words and concepts that have
gained popularity, such as cloud computing and mobility. As a consulting
firm, our clients are increasingly requesting these features to enhance
their processes and handle larger volumes of data. This includes
leveraging analytics and exploring the Internet of Things.

\begin{figure}[!h]
    \centering
    \includegraphics[page=14, trim = 2cm 5cm 2cm 0cm, clip, width=\textwidth]{images/01 - KPMG_IT Advisory Services_ENG_v1.pdf}
\end{figure}

These are the keywords that define the evolution of the
market: accessibility, collaboration, and synchronization are the
driving forces behind this shift. One of the leading software vendors in
this field is SAP, an international company that has been instrumental
in shaping this framework. Until about five years ago, most companies
were developing and implementing their systems on-premise. This meant
that they purchased licenses and installed the software on their own
servers.

Nowadays, many vendors, including SAP, are adopting a philosophy of
moving everything to the cloud. This means that companies looking to
implement SAP solutions no longer need to install them on their own
servers. Instead, everything is accessible through the cloud. The core
business processes, such as finance and accounting, can still be kept
on-premises using the core ERP system. However, specific processes, like
purchasing, can be linked to dedicated cloud solutions.

For instance, consider the purchasing processes of creating a purchase request, issuing a purchase order, and receiving goods. These processes can be efficiently managed using a dedicated cloud solution called Ariba\footnotemark{}, which is tightly integrated with the core ERP system. This is the direction the market is moving towards, and it is important to consider this when making a business case.

In summary, the trend is to move towards cloud-based solutions, with
core processes remaining on-premises and specific processes being
managed through dedicated cloud solutions.

\footnotetext{Ariba is a cloud-based procurement management solution provided by SAP. It integrates with the core ERP system to automate complex workflows, making it easier for employees to search for goods and services, collaborate with suppliers, and manage approvals and invoices \href{https://learning.sap.com/products/intelligent-spend-management/ariba/procurement}{Source 1}.}


\subsection{Integration Challenges and
    Strategies}\label{integration-challenges-and-strategies}

\subsubsection{Choice of an ERP}
\begin{figure}[!h]
    \centering
    \includegraphics[page=22, trim = 2cm 3cm 2cm 0cm, clip, width=\textwidth]{images/01 - KPMG_IT Advisory Services_ENG_v1.pdf}
\end{figure}

\begin{figure}[!h]
    \centering
    \includegraphics[page=23, trim = 1cm 3cm 1cm 0cm, clip, width=\textwidth]{images/01 - KPMG_IT Advisory Services_ENG_v1.pdf}
\end{figure}

In terms of the main software vendor for the ERP there are Oracle and SAP, but first there are some keywords that are important to consider. The structure of the
software is based on modules, with each module representing a specific
area, department, or process. For example, there are modules for the
financial department, controlling, sales, purchasing, and production.
Each module is connected to the others, meaning that actions in one module
can impact other modules. For instance, creating a purchase requisition
or order in the purchasing module may require entering financial objects
that are key elements in the controlling or finance modules.

\subsubsection{ERP Projects}

\begin{figure}[!h]
    \centering
    \includegraphics[page=38, trim = 2cm 2.5cm 2cm 0cm, clip, width=\textwidth]{images/01 - KPMG_IT Advisory Services_ENG_v1.pdf}
\end{figure}

Implementing an ERP system involves three key dimensions: technology,
processes, and people. There are several risks associated with ERP
projects, as shown by the results of previous campaigns. Project
management is crucial, as a poorly structured project plan can lead to
issues. Setting clear business objectives from the beginning is
essential to manage expectations and measure the benefits of the
implementation. Managing complexity and making the right software
selection are also important factors.

When starting a project, it is necessary to collect functional
requirements by interviewing process owners and key users from each
department. These requirements need to be translated into the language
of the ERP system to proceed with the implementation. The implementation
process includes system setup, testing in a quality environment, user
acceptance testing, and training for all users in the company.

\begin{figure}[!h]
    \centering
    \includegraphics[page=39, trim = 2cm 4cm 2cm 0cm, clip, width=\textwidth]{images/01 - KPMG_IT Advisory Services_ENG_v1.pdf}
\end{figure}

Choosing the right software and implementing it can be complex.
Consulting companies often have their own approach and tools to
accelerate the process. It is important to consider the business needs
and plan for the future when selecting and implementing an ERP system.

\begin{figure}[!h]
    \centering
    \includegraphics[page=48, trim = 2cm 1.5cm 2cm 0cm, clip, width=\textwidth]{images/01 - KPMG_IT Advisory Services_ENG_v1.pdf}
\end{figure}

\begin{figure}[!h]
    \centering
    \includegraphics[page=53, trim = 1cm 1.5cm 1cm 0cm, clip, width=\textwidth]{images/01 - KPMG_IT Advisory Services_ENG_v1.pdf}
\end{figure}

In terms of transformation projects, change management is crucial. It
involves explaining the benefits and changes to the organization and its
departments. This includes cost savings and potential revenue increases.
It is important to consider the current state of the legacy system and
its ability to manage processes. If acquiring another company,
integration planning is necessary, whether it involves adopting their
system or finding a common system. The selection of software depends on
specific requirements and constraints, such as supporting collaboration
with suppliers. The relationship between processes, technology, and
change management is essential. In the case of a global company,
localization requirements and multiple instances may need to be
considered. The integration of acquired companies and their IT landscape
is also a strategic consideration.
% \section{KPMG on Low Code}\label{introduction-to-low-code-development}

I'm not sure which one of you is familiar with this new term, but it's
gaining a lot of attention globally. Although it may not be as popular
in Europe and Italy, there is another approach to software development
that is worth exploring. While technologies like that may be the current
trend, they are not the only ones that software engineers should focus
on in the next decade. There are new methods emerging that make software
development easier, faster, and more secure. Now, let's dive into the
details of low-code development, which is the topic we will be
discussing today.

\subsection{Definition and Scope}\label{definition-and-scope}

\begin{figure}[!h]
    \centering
    \includegraphics[page=7, trim = 2cm 2cm 2cm 0cm, clip, width=\textwidth]{images/02 - KPMG_intelligent_automation_2.pdf}
\end{figure}

Low-code development is more than just a platform or a cloud solution.
It encompasses a combination of these elements. Now, let's move on
to discuss the major technologies involved in low-code development.

\subsection{Historical Context and
    Evolution}\label{historical-context-and-evolution}

Over the past decade, companies have
been increasingly focused on optimizing and streamlining their core
processes, such as procurement and operations. However, it's important
for companies to consider not only their core operations but also the
surrounding components that impact their overall efficiency. This
includes visualizing and automating common user tasks and leveraging
data to improve communication and collaboration with other departments.

In order to foster innovation, companies need to shift their attention
and resources towards individuals who can generate new ideas and
solutions. One common approach is to automate repetitive tasks or
non-value-added activities. This trend started years ago in large US
companies, including industrial giants like GM and banks, as well as
government institutions. Initially, the focus was on workforce
optimization rather than investment analysis. This was a simple and
straightforward solution at the time.

\subsection{Global Trends and
    Adoption}\label{global-trends-and-adoption}

In order to maintain our jobs and avoid outsourcing to
low-income countries like India, we need to take action. We can achieve
this by finding skilled individuals in these countries who can perform
our job effectively and at a lower cost. By doing so, we can reduce our
expenses and present a compelling case to management. This strategy is
commonly referred to as ``Outsourcing'' It is crucial for all
of us in the driving industry to understand and embrace this concept.

\subsection{Automation and Low-Code}\label{automation-and-low-code}

In these companies, there is a strong focus on automating activities
that were previously done by humans. This is the first aspect of
low-code development: automation. The idea behind this approach is to
find ways to automate these activities using robotic process automation
software. This technology is not new; it has been around for 14 years
and is commonly used in various applications. By implementing this
technology, companies can streamline their processes and reduce costs.

\subsection{Analytics and Data}\label{analytics-and-data}

The second aspect we will discuss is analytics. Analytics involves the
ability to analyze data collected from various sources within an
enterprise software system. It is important to create reports that not
only provide valuable insights but also optimize cost-effectiveness.
Analytics plays a crucial role in boosting automation. Automation relies
on both software and data, similar to how humans require data to perform
tasks efficiently. For instance, if a human needs to transfer data from
one window to another, they need the data to be organized and in a
suitable format. This ensures greater efficiency and effectiveness in
the overall process.

\subsection{Artificial Intelligence and
    Low-Code}\label{artificial-intelligence-and-low-code}

Lastly, but certainly not least, is the
application. Sometimes, you need tools to help you understand data and
make informed decisions, especially in automated processes. This is
where artificial intelligence (AI) comes into play. When AI intersects
with low-code development and automation, it creates a powerful
combination that can drive international competition.

\subsection{Low-Code as a Competitive
    Advantage}\label{low-code-as-a-competitive-advantage}

In the realm of international competition, the convergence of three key
elements---software, robotics, and business operations---is of utmost
importance. This combination allows for the creation of applications and
technologies that can significantly enhance business operations. One
notable advantage of this technology is its ability to automate crucial
tasks, resulting in increased efficiency. Moreover, the speed at which
these applications can be developed is truly remarkable. With this
technology, it is possible to create and deploy applications rapidly.

\subsection{The Future of Software
    Development}\label{the-future-of-software-development}

In the future, the landscape of software development will be drastically
different. Currently, companies rely on a large number of employees to
write code and build applications. However, with the advancements in
analytics, artificial intelligence, and automation, this process will be
revolutionized. These technologies will replace the need for manual
coding, making it much faster and more efficient.

As a result, the role of universities in teaching high-level programming
languages may become less important. The traditional approach to
software development will be replaced by low-code development, which
allows for faster and more streamlined application building. This means
that individuals who are currently studying programming may need to
reconsider their career paths, as the demand for manual coding skills
may decrease significantly in the next century.

It's important to embrace these changes and adapt to the future of
software development. By understanding the potential of low-code
development and the power of analytics, AI, and automation, individuals
can position themselves for success in this evolving industry.

\subsection{Low-Code Market Value}\label{low-code-market-value}

Why am I saying this? Because even today, if I want to develop a simple
enterprise application that requires some coding, I would
have needed to hire a Java or .NET programmer in the past. It would have
taken them around 10 to 20 days to complete the task. However, with
low-code development, I can do it myself in just one day. This
demonstrates the efficiency and time-saving benefits of low-code.

In the near future, I may not even need a person to configure the
low-code platform. Companies like Microsoft are continuously improving
their low-code tools, making them more user-friendly and accessible.
This evolution in low-code development is not just about individuals,
but it also has a significant impact on industries. Low-code has already
revolutionized billion-dollar infrastructures, not just small
businesses.

The global market value of low-code is currently around \$100 billion
per year, indicating a significant shift away from traditional
development methods.

\section{The Impact of Low-Code}\label{the-impact-of-low-code}

\subsection{Job Market and Industry
    Changes}\label{job-market-and-industry-changes}

My role here is to acknowledge a significant shift in the job market. In
my interviews with colleagues, I've encountered individuals who aspire
to be full-stack programmers, back-end developers, and content
developers. While these roles are valuable, it's important to recognize
that they may become less prevalent in the future. Companies like ours, have chosen to focus on the strength of technology rather
than relying on a large number of developers spread across the globe. We
have witnessed the growing importance and popularity of technologies
like cybersecurity and bandwidth growth. This is why we have made the
decision to invest in these technologies and share this information with
you.

\subsection{The Role of Traditional Software
    Development}\label{the-role-of-traditional-software-development}

I understand your concerns about your ability to excel in your job and
the challenges you face in the software industry. It can be difficult
when there is a lack of resources and support. However, it's important
to remember that the industry is constantly evolving, and there are new
opportunities emerging.

One such opportunity is low-code development, which allows individuals
without extensive coding experience to create applications. Low-code
platforms provide a visual interface and pre-built components that
simplify the development process. This means that even if you don't
consider yourself a skilled engineer, you can still contribute to
software development.

The rise of low-code development has also led to increased automation in
the industry. Tasks that were once time-consuming and manual can now be
automated, freeing up valuable time for more complex work. This
automation is particularly beneficial in the front-end development
space, where there is often a shortage of skilled individuals.

However, it's important to acknowledge that the traditional software
development industry is still relevant and necessary. While low-code
development offers advantages, there are still complex projects that
require the expertise of experienced engineers. It's about finding the
right balance and leveraging the strengths of both approaches.

In conclusion, the software industry is evolving, and low-code
development is becoming a competitive advantage. It provides
opportunities for individuals without extensive coding backgrounds and
enables automation in the development process. However, traditional
software development still plays a crucial role in tackling complex
projects. It's important to adapt to these changes and embrace new
technologies to stay relevant in the industry.

\subsection{New Opportunities in Software
    Industry}\label{new-opportunities-in-software-industry}

There is a new and rapidly growing software industry that is gaining
momentum. It may seem like a small trend, but it has actually been
developing for the past 20 to 30 years. Initially, there were only a few
companies involved, but now there are countless players in the market.
These companies are often overshadowed by larger, more well-known
organizations, but they are making significant contributions in various
sectors, including healthcare and media.

The emergence of this new software industry has created numerous job
opportunities, but there is a shortage of skilled professionals who are
familiar with the technology. One reason for this shortage is that many
people are not yet aware of the potential of this industry.
Additionally, traditional software developers may be hesitant to
transition to this new field.

However, it is important to recognize the potential and embrace the
changes that this industry brings. It offers exciting prospects and
opens up new avenues for innovation. It's time to break away from
traditional norms and explore the possibilities that this new software
industry has to offer.

\section{Low-Code in Practice}\label{low-code-in-practice}

\subsection{Intelligent Automation}\label{intelligent-automation}

So, returning to our presentation, these are the tasks that
someone working in intelligent automation can perform. It's not just
about developing solutions; it involves IT professionals engaging in a
wider range of activities. It's not like being a software developer,
where you gather requirements, develop an application, and then test and
deploy it. In intelligent automation, we also have the task of
discovery. We engage in value-added activities to understand how
automation can be applied to companies, build those automation
solutions, and incorporate elements of cognitive and artificial
intelligence. And, ultimately, we also handle management. Let me provide
an example. One of our projects serves as a very interesting use case.

\subsection{Use Cases and Efficiency}\label{use-cases-and-efficiency}

For an automotive company that increased its production of luxury cars,
there was a challenge related to checking the labels on each vehicle.
These labels contain important information, such as the manufacturing
location and pressure specifications, and must comply with local
regulations and be in the language of the market where the car is sold.
Previously, employees on the production line manually checked each label
for accuracy and compliance.

To address this issue, a discovery process was conducted to identify
opportunities for automation. A solution was developed using a
combination of low-code development, cognitive technology, and
artificial intelligence. An app was created for tablets used by the
employees on the production line. They could simply take a picture of
the labels on the car doors and other locations, and an Amazon Web
Service configured with label recognition capabilities would
automatically analyze the image. The system would determine if the label
was correct, in the right position, and compliant with regulations. The
results were recorded and archived in the tool.

By implementing this automated solution, the automotive company was able
to reduce the number of employees needed for label checking from five to
just a couple. This resulted in significant cost savings for the
company.

\subsection{Tools and Technologies}\label{tools-and-technologies}

At our company, we go beyond just being software engineers. We strive to
find the best solutions and automation techniques to add value to our
clients' businesses. We use various tools, such as Microsoft, AWS, and
others, to create high-level architectures for automating human
activities. One of the tools we use is Robotic Process Automation (RPA),
which allows us to automate repetitive tasks. We also utilize HR and
cognitive services, as well as low-code development, to build
applications. Think of us as painters who combine different technologies
to automate work processes. This approach is incredibly powerful and
transformative.

\subsection{Traditional vs.~Low-Code
    Development}\label{traditional-vs.-low-code-development}

If we look back 10 or 20 years ago, the only option for software
development was writing code in specific programming languages like .NET
or C\#. However, today we have a wider range of technologies
available to us. This has led to the emergence of low-code development.

Traditional development is a manual process carried out by humans, which
can be prone to errors and requires significant resources. Developers
often spend a lot of time searching for code, trying to understand it,
and sometimes making mistakes that require further investigation. On the
other hand, low-code development offers a different approach.

With low-code, developers have access to a palette of pre-built
components and services that can be reused. Many of these components
have already been created by other developers, allowing for efficient
reuse of their work. This is the essence of low-code development: the
ability to create applications without having to write traditional code.

\section{The Future of Low-Code}\label{the-future-of-low-code}

\subsection{Predictions and Market
    Trends}\label{predictions-and-market-trends}

Low-code development is gaining popularity due to its ease of use and
ability to bridge the gap between business needs and software
development. Compared to traditional software development, low-code
allows for faster results and the ability to showcase them more quickly.
This is particularly important because by 2025, the majority of
enterprise application development is expected to be done using
low-code. This shift towards low-code presents a significant opportunity
for industries to adopt this approach. The low-code market is projected
to reach \$30 billion by 2025, making it crucial for businesses to
embrace this trend and stay competitive.

\subsection{Digital Solutions and Legacy
    Modernization}\label{digital-solutions-and-legacy-modernization}

Low-code development offers a range of possibilities for various
applications. One key benefit is the ability to create new digital
solutions that can be used both internally and externally. These
solutions can be utilized by employees within the company as well as by
consumers who want to interact with the organization.

Another valuable use of low-code is the modernization of legacy
applications. Many industries, such as banking, still rely on outdated
applications developed in the 70s or 80s. With low-code, companies can
transition these old technologies to more modern platforms by developing
low-code applications.

Additionally, low-code can be used to streamline and consolidate
applications across different departments. In the past, companies would
build separate applications for each department, resulting in a lack of
cohesion. Low-code allows for the unification of these applications into
a single, integrated system.

Lastly, low-code enables the automation and orchestration of existing
processes, as well as the design of new processes. This means that
companies can automate repetitive tasks and workflows, improving
efficiency and productivity.

These four areas---creating new digital solutions, modernizing legacy
applications, unifying applications across departments, and automating
processes---are the primary use cases for low-code development.

\subsection{Low-Code Landscape}\label{low-code-landscape}

\begin{figure}[!h]
    \centering
    \includegraphics[page=14, trim = 2cm 2cm 2cm 0cm, clip, width=\textwidth]{images/02 - KPMG_intelligent_automation_2.pdf}
\end{figure}

\begin{figure}[!h]
    \centering
    \includegraphics[page=15, trim = 2cm 2cm 2cm 0cm, clip, width=\textwidth]{images/02 - KPMG_intelligent_automation_2.pdf}
\end{figure}

In the landscape of low-code development, there are different types of
applications that can be built. One example is multi-experience
platforms, which allow you to develop applications that can be deployed
on mobile devices, tablets, or web browsers. These platforms focus on
providing a seamless user experience across different devices.

Another type of platform is functional extension platforms. These
platforms are used to enhance existing enterprise software by
configuring and augmenting processes. They enable businesses to extend
the functionality of their software without the need for extensive
coding.

For those who are new to coding or prefer a more user-friendly approach,
there are citizen developer platforms. These platforms require no coding
knowledge and are designed to be easy to use for anyone, including
non-technical individuals. They are often used by companies to create
small-scale applications quickly and efficiently.

Lastly, there are business process management platforms. These platforms
are used to digitize complex processes within organizations. By
leveraging these technologies, businesses can develop solutions that
streamline and automate intricate processes.

Overall, the low-code landscape offers a range of platforms tailored to
different application development needs, from multi-experience platforms
to functional extension platforms, citizen developer platforms, and
business process management platforms.

\section{Hands-On Exercise}\label{hands-on-exercise}

\subsection{Design Activity and User
    Experience}\label{design-activity-and-user-experience}

Low-code development can be used in various scenarios, such as digital
transformation initiatives, workflow automation, creating specialized
software for niche industries where enterprise software is not
available, and modernizing legacy systems. The versatility of low-code
allows for a wide range of applications.

For today's exercise, we will focus on a low-code platform and simulate
a project similar to what we do in real-life low-code projects. One of
the initial phases in a low-code project is the design activity. During
this phase, we need to consider the process, the data involved, the
users of the application, and the user interface.

Designing a user interface is a complex task in today's world. We need
to develop applications that are easy to use for everyone and accessible
across multiple devices. Low-code software can assist in this process by
providing a palette of optimized graphic layouts for mobile, laptop, and
public use. However, creating a user interface that meets these
requirements remains challenging.

For today's exercise, we will focus on user experience (UX) and user
interface design (UI).

\subsection{ACME Company Case Study}\label{acme-company-case-study}

\begin{figure}[!h]
    \centering
    \includegraphics[page=17, trim = 2cm 2cm 2cm 0cm, clip, width=\textwidth]{images/02 - KPMG_intelligent_automation_2.pdf}
\end{figure}

ACME, a company, has approached us to develop a workforce time report. A
time report is essentially a calendar where individuals working on
projects or for a company can log their time. They record the number of
hours spent on specific tasks, such as training, project work, or email
responses. ACME has requested us to design the entire solution,
including the user interface.

When it comes to the user interface, our usual approach is to create a
wireframe. A wireframe is a sketch of the user experience that takes
into account the technical requirements of IT software. It needs to be
compatible with webpages, easy to use, and mobile-friendly. Designing a
user experience is a complex task that requires careful consideration.

To facilitate this process, we recommend using a tool like \href{https://www.drawio.com}{Draw.io}. It
allows us to create detailed wireframes that align with ACME's needs and
preferences. By collaborating as a group, we can ensure that the final
design meets all the necessary criteria for a successful workforce time
report.

\subsection{Success Criteria and Business
    Rules}\label{success-criteria-and-business-rules}

One of the free tools available on the web for building a user
experience is wireframing. While you can use PowerPoint, I strongly
recommend using a tool specifically designed for this purpose. To create
the wireframes for ACME Company, we need to start with an on-page for
the employee. As an employee, I would need to enter my username and
password, and there could also be additional authorization steps such as
multi-factor authentication or a one-time password. We also want to
include navigation to the time report page and an expense report page.

The expense report page is where employees can enter their expenses,
such as travel expenses. This page should allow users to input details
about the expense, including the date, type, receipt upload, amount, and
location. It's important to consider all the necessary information for
expense reporting.

When designing the wireframes, it's crucial to keep in mind the success
criteria. As consultants, we are always evaluated based on these
criteria. One important aspect is the navigation between pages. It's
common for teams to focus on designing beautiful pages but overlook the
seamless transition between them. Users should be able to easily
navigate back and forth between pages, as mistakes can happen during
data entry.

Another consideration is the user interface (UI). Avoid using large
white boxes on a page, as this can be a common mistake. Just like you as
a consumer wouldn't appreciate a lackluster UI on your iPhone or
Android, it's important to judge your own design from a user's
perspective. Additionally, it's essential to define the business rules
behind each page. These rules will guide the functionality and behavior
of the application.

By taking these factors into account and creating wireframes that
prioritize user experience, navigation, and a visually appealing UI, we
can ensure the success of the ACME Company application.

\section{Closing Remarks and Exercise
  Instructions}\label{closing-remarks-and-exercise-instructions}

For example, let's consider a validation process. It's important to
ensure that expenses are only entered for the days worked. Keep in mind
that these validations are not optional---they are necessary.
Additionally, it's crucial to make sure that the user interface (UI) is
responsive and adaptable to different devices. Instead of relying on
adjusting the power pointer, it's recommended to utilize the UI, which
is accessible across all devices. Remember, we are not just focusing on
time, but also on expenses. It's important to consider the user
experience and not overlook the needs of the users. Lastly, building
security is essential. An authentication mechanism should be implemented
to protect sensitive data. This exercise is designed to help you
understand the group's requirements. Since there is a lot of work to be
done and a meeting scheduled for later this afternoon, it would be
beneficial for you to start working on it now.

\section{Web Information Systems - I Lecture}

\subsection{Introduction and Definition of WIS}\label{introduction-to-web-information-systems}

Good morning! In this class, we will be discussing web information
systems. As usual, you can find the slides on WEBEEP. Let's begin by
defining what a web information system is.

\begin{figure}[!h]
  \centering
  \includegraphics[page=2, trim = 2cm 3cm 2cm 4.5cm, clip, width=\textwidth]{images/03 - Web_Information_Systems.pdf}
\end{figure}

A web information system is a system that facilitates communication
between machines using either the public internet or an IP-based private
VPN. Users can access the system's functionalities through a web
browser. This definition provides a broad understanding of what a web
information system entails.

\paragraph*{Defining Web Information Systems}
The public internet and the IP protocol, along with browsers, have
brought about significant changes in information systems. Any system
that meets two criteria can be considered a web information system.
Firstly, it must be accessible through a browser. Secondly, it should be
connected to the private VPN of the company. This means that even a
traditional ERP with its core functionalities can be classified as a web
information system if it meets these criteria.


\subsection{Impact of Technical Innovations}
These two technical
innovations---internet and browsers (user interfaces)---have resulted in a transformation of the technologies
utilized, even by components of web information systems that were
created prior to the advent of the web. In essence, ERP providers have
reimagined the interface and core technologies of their packages to
align with this new paradigm.

\subsubsection{The Internet and Its
  Influence}\label{the-internet-and-its-influence}

\begin{figure}[!h]
  \centering
  \includegraphics[page=4, trim = 2cm 2.5cm 2cm 1cm, clip, width=\textwidth]{images/03 - Web_Information_Systems.pdf}
\end{figure}

Now, let's discuss the main innovation that revolutionized web
information systems: the internet. The internet connected companies to
individual customers for the first time, marking a significant change in
the business landscape. This not only impacted individual customers but
also the employees working for these companies. The ability to connect
the company's potential market, especially retail customers, through a
public network was seen as a monumental shift.

During the late 90s, when the internet was introduced, there was a great
deal of hope and belief in the transformative power of this new
technology. There was a hype surrounding the web, with expectations that
it would drastically change people's lives in a short period of time.
The paradigm of exponential organization emerged, where companies
without the burden of physical shops or channels to reach customers
could experience rapid growth. As a result, the value of web-based
companies experienced significant growth around the year 2000.

\paragraph{The Dot Com Bubble}\label{the-dot-com-bubble}

The dot com bubble refers to a period when the value of stocks of
companies with websites ending in ``.com'' skyrocketed, far exceeding
their actual revenues and assets. However, this rapid growth was
eventually followed by a significant stock market crash, one of the
largest in recent history. This sudden shift from extreme optimism to
pessimism marked a turning point in the market.

\paragraph{Web Technology in Traditional
  Companies}\label{web-technology-in-traditional-companies}

The web has undoubtedly revolutionized processes for both traditional
and new companies, creating new industries while also causing the
decline of others. However, the pace at which society adapted to these
changes was slower than the expectations of the stock exchange. For
traditional companies to fully benefit from web innovation, they must
first go through the implementation and integration of their information
systems, even with traditional technologies. Without this, it is
challenging for them to compete effectively on the web, unless they have
a history of being highly innovative in their approach to information
systems and are prepared to utilize the web as a new channel to reach
customers.

For instance, companies that lack a unified data system, which is a
fundamental aspect of the ERP paradigm, will struggle to provide a
satisfactory web service. The web is a more objective platform compared
to a physical store, as there is no personal interaction. Customers have
certain expectations for service quality, which are set by industry
leaders. Consequently, the web can actually make the customer
relationship more challenging rather than easier. Only companies that
have successfully integrated previous technologies into their operations
have been able to fully exploit the potential of the web.

\paragraph{Extended ERP and E-Commerce}\label{extended-erp-and-e-commerce}

When examining the functional architecture of ERP systems, it becomes
clear that extended ERP is more influenced by the web than core ERP. The
web has opened up new opportunities, particularly in e-commerce, which
continues to evolve. Extended ERP, which is enabled by the web,
encompasses various functionalities that must be supported by a
well-integrated and IT-supported internal process.

One of the key enabling technologies for extended ERP is CRM, with the
web playing a significant role. While companies have been using
technology, such as Salesforce automation, to interact with customers
for some time, the web has made it easier to showcase services, prices,
and conditions, as well as compare alternative suppliers. This has
created a need for integration across different channels, as any
mistakes or inaccuracies on the web are visible to all. It is no longer
acceptable for customers to have to visit a retail shop for information;
it must be readily available and accurate on the web.

\subsection{E-Commerce and Its
  Evolution}\label{e-commerce-and-its-evolution}

Companies have recognized the importance of integrating customer
information into a single database, commonly known as a customer
database or CRM. This integration, whether referred to as omni-channel
or multi-channel, has become increasingly important as the web has
become the primary access point for both customers and internal users.

\begin{figure}[!h]
  \centering
  \includegraphics[page=5, trim = 2cm 5cm 2cm 4.5cm, clip, width=\textwidth]{images/03 - Web_Information_Systems.pdf}
\end{figure}

Initially, the web was seen as a platform for e-commerce, allowing
companies to sell their products online and reach a wider market. This
was seen as a significant opportunity. E-commerce refers to the buying
and selling of products and services online, primarily targeting retail
customers. On the other hand, when referring to business customers, the
term e-business is used. For example, if a company purchases from
another company online, it is considered e-business. Similarly, if a
public administration buys from a supplier online or provides services
to other public institutions or citizens, it is referred to as
e-government.

In summary, e-commerce focuses on reaching the market and selling
products online, while e-business and e-government encompass a broader
range of online transactions involving business customers and public
administrations, respectively.

\subsection{The Dot Com Trend and Its
  Aftermath}\label{the-dot-com-trend-and-its-aftermath}

In the early days of the dot com trend, many companies initially set up
their e-commerce operations as separate units or even separate
companies. This was a common approach when a new technology brought
about significant changes. These companies lacked the necessary
expertise to effectively manage and leverage the new technology, so they
sought the help of consultants and formed internal teams to handle it.
Over time, these teams evolved into dedicated units focused on the new
technology, such as the web.

\begin{figure}[!h]
  \centering
  \includegraphics[page=6, trim = 1.5cm 4cm 2cm 3cm, clip, width=\textwidth]{images/03 - Web_Information_Systems.pdf}
\end{figure}

This pattern of innovation is not unique to the web; it has been
observed in various industries. However, during the dot com era,
companies that were perceived as pure dot coms, operating solely on the
web without any traditional brick-and-mortar presence, were highly
valued in the stock market. This led even traditional companies to
establish separate units or even separate companies with their own
distinct brands to tap into the potential of the web. For instance, in
2000, Deutsche Bank launched Bank 24, a web-only bank operating as a dot
com.

However, this trend came to a halt after the dot com bubble burst and
the subsequent stock market failure. Even Deutsche Bank eventually
integrated Bank 24 back into its core services, recognizing the need to
reassess their approach in light of the challenges faced by dot com
companies.

\subsection{Integrating Web and Traditional
  Channels}\label{integrating-web-and-traditional-channels}

Companies have recognized the importance of integrating the web as an
additional channel for their existing customers. There are a few reasons
behind this realization. Firstly, companies understood that relying
solely on the web for growth and stock market success was no longer
feasible after the dot com bubble burst. Secondly, they realized that
web-based operations still require physical processes and the management
of people, offices, and complementary services. This led to the
recognition of the synergies between web-based companies and traditional
structures, as Deutsche Bank recognized.

To leverage these synergies, companies brought web-based operations back
into their traditional structures and began integrating the web as a
channel alongside other existing channels. This integration process was
facilitated by customer relationship management (CRM) practices.
Companies started by creating a shared customer database and analyzing
it to understand the evolving habits of web customers. Based on these
insights, they developed innovative services to cater to the changing
needs of web customers over time.

\begin{figure}[!h]
  \centering
  \includegraphics[page=7, trim = 2cm 1.5cm 2cm 1.5cm, clip, width=\textwidth]{images/03 - Web_Information_Systems.pdf}
\end{figure}

Initially, customers were skeptical about making purchases on the web,
especially for high-value or high-quality products. They preferred
traditional brick-and-mortar shops. However, over time, customer
behavior shifted, and e-commerce gained traction across various
industries. The pace of this shift varied across industries, with some
experiencing faster adoption than others. Nevertheless, e-commerce has
become prevalent in almost all industries, including fashion, grocery,
tourism, banking, trading, and learning. Each industry has its own
specific terms and practices related to e-commerce.

\subsection{The Role of E-Commerce
  Post-COVID-19}\label{the-role-of-e-commerce-post-covid-19}

When the internet was first introduced, there were predictions about
which industries would be most impacted by e-commerce and how quickly it
would replace traditional channels. However, it took nearly 20 years for
e-commerce to become pervasive, serving only a small percentage of
transactions and clients. The COVID-19 pandemic changed this
dramatically, giving e-commerce a significant boost. It quickly
developed and penetrated the market, reaching a larger audience. While
there has been a slight decline after the initial surge, e-commerce has
still experienced substantial growth and increased penetration.

\subsection{Customer Journey and Online
  Information}\label{customer-journey-and-online-information}

To illustrate this point, let's consider the example of grocery
shopping. Despite the pandemic, only a small percentage of people
actually do their grocery shopping online. This indicates that the
experience of shopping in a traditional store or through traditional
distribution methods is still considered superior to the online
experience. As a result, the current trend is to view online shopping as
complementary to other channels, rather than a standalone option.

\begin{figure}[!h]
  \centering
  \includegraphics[page=8, trim = 1.5cm 3cm 2cm 3cm, clip, width=\textwidth]{images/03 - Web_Information_Systems.pdf}
\end{figure}

When customers engage with a company, they go through a journey that
involves multiple channels at different stages of the transaction. It
all begins with a need, which may arise from seeing an advertisement or
simply having a requirement. The first step in this journey is typically
going online to gather information. This initial step is crucial in the
transaction process and is usually done through the web. When customers
are unsure of where to make a purchase, they turn to the internet to
find information. This helps them make decisions about the provider, the
brand, and the specific product they want to buy. With the advent of
social media, customers also rely on product reviews to guide their
choices. This search phase of the transaction, where customers seek
information, accounts for approximately 90\% of visits to a company's
website.

\subsection{Designing Effective Web
  Information}\label{designing-effective-web-information}

The main purpose of websites is to provide information to customers who
are looking for products or learning about a company for the first time.
Typically, websites include a general presentation of the company, its
mission, organizational structure, job opportunities, news and events,
and information about its products. Some older websites may offer a
downloadable product catalog in PDF format or a multimedia online
catalog where customers can browse through the products. These websites
often provide recommendations for related products that are commonly
purchased together. Contact information, such as the call center,
company size, and location, is also provided, sometimes with a map for
easy navigation.

When designing the information services of their website, companies are
primarily concerned with providing a positive first impression to
customers who have no prior knowledge of the company. This can be
challenging for companies, as they need to carefully design the
website's navigation structure. It is often best to follow a navigation
structure that is similar to what competitors use, as customers have
certain expectations about where to find specific information on a
company's homepage. Keeping the information provided on the website
standardized is a good idea, although it should be regularly updated as
company information, including the mission, can change over time. To
update the information on their website, companies need to gather the
necessary information from various stakeholders within the organization
who possess the relevant knowledge. This can be a complex task, as
information is often scattered across different departments or
individuals. The website serves as a central access point for customers
to find all the information they need.

\subsection{Quality Criteria for E-Commerce
  Sites}\label{quality-criteria-for-e-commerce-sites}


\begin{figure}[!h]
  \centering
  \includegraphics[page=9, trim = 1.5cm 4.5cm 3cm 4cm, clip, width=\textwidth]{images/03 - Web_Information_Systems.pdf}
\end{figure}

In order to provide consistent information to customers on the web,
information owners must cooperate. There are two approaches to
collecting this information. The first approach is federated, where
there is one central site with a thin layer of information and multiple
local sites where individual business units can provide their own
information. The landing page serves as a collection of links to the
different business units' websites. While this approach allows for
information retrieval to be closer to the business units, it often
results in a website that lacks standardization and has lower overall
quality.

The second approach is to establish an editorial committee that
constantly updates the information on the website. This committee seeks
information from the different business units and follows a common
template for all units, ensuring standardization. Although this approach
is more expensive due to the need for an editorial committee and
agreement on the information to be provided, it generally leads to
higher quality.

\begin{figure}[!h]
  \centering
  \includegraphics[page=10, trim = 2cm 2cm 2cm 4cm, clip, width=\textwidth]{images/03 - Web_Information_Systems.pdf}
\end{figure}

When considering the quality criteria for an e-commerce site, there are
several factors to consider. First, customers expect the content to be
complete, dependable, and correct. They also want consistency in the
information provided by the call center. Second, the structure of the
website should allow customers to easily navigate and understand the
information. A centralized structure is preferable in this regard. The
presentation of the website should have a coherent graphic style to
instill trust and confidence in the company. The layout and position of
information should be intuitive and standardized. Navigation should also
be intuitive, with multiple paths to reach the same information. The
website should be well-connected, allowing customers to navigate between
pages without having to return to the homepage. The ease of navigation,
including the ability to go back multiple steps, is crucial.

In conclusion, the quality criteria for an e-commerce site include
content completeness and dependability, a clear and intuitive structure,
a coherent graphic style, and easy navigation with multiple paths to
reach information.

\section{Web Information Systems - II Lecture}
\subsection{Introduction to Web Information Systems - Part
  Two}\label{introduction-to-web-information-systems---part-two}

Welcome to part two of our series on web information systems. In this
installment, we will delve deeper into the fascinating world of
web-based information systems and explore their various components and
functionalities. So, let's jump right in and continue our exploration of
this exciting topic.

In part one, we introduced the concept of web information systems and
discussed their importance in today's digital age. We learned that web
information systems are designed to collect, process, store, and
disseminate information over the internet. They play a crucial role in
facilitating communication, collaboration, and decision-making in
various domains, such as business, education, healthcare, and
government.

Now, let's move on to the key components of web information systems. At
the heart of any web information system is the web server, which acts as
the central hub for storing and serving web pages and other resources.
The web server receives requests from clients, such as web browsers, and
responds by sending the requested information back to the client.

Another essential component of web information systems is the database.
The database stores and organizes the vast amount of data that is
generated and used by the system. It allows for efficient storage,
retrieval, and manipulation of data, ensuring that the system can handle
large volumes of information effectively.

Web applications are another critical element of web information
systems. These applications are responsible for processing user
requests, generating dynamic content, and providing interactive
features. They enable users to perform various tasks, such as submitting
forms, conducting searches, and accessing personalized information.

To ensure the security and integrity of web information systems,
authentication and authorization mechanisms are implemented.
Authentication verifies the identity of users, while authorization
determines what actions they are allowed to perform within the system.
These mechanisms help protect sensitive information and prevent
unauthorized access.

In addition to these components, web information systems often
incorporate other technologies and tools, such as content management
systems, search engines, and analytics platforms. These tools enhance
the functionality and usability of the system, allowing for efficient
content creation, search capabilities, and data analysis.

In conclusion, web information systems are complex and multifaceted
systems that enable the collection, processing, storage, and
dissemination of information over the internet. They consist of various
components, including web servers, databases, web applications,
authentication and authorization mechanisms, and additional tools and
technologies. Understanding these components is crucial for building and
maintaining robust and effective web information systems. Stay tuned for
part three, where we will explore the design and development process of
web information systems.


\subsection{Gaining Online Visibility}\label{gaining-online-visibility}

Now that we have discussed the quality criteria for websites, let's take
a closer look at what makes a website truly exceptional.

\subsubsection{Search Engines and User
  Behavior}\label{search-engines-and-user-behavior}

Despite having a high-quality website, companies still need to ensure
that their online presence is visible. So, how do companies initiate the
search phase of a transaction? Well, they start by going online and
conducting a search on popular search engines like Google or Bing. It's
worth noting that Google is the most widely used search engine, with a
near monopoly in certain countries, particularly when it comes to web
searches. When conducting a search, people usually look for a specific
product or service.

\subsubsection{Google's Dominance and Search
  Patterns}\label{googles-dominance-and-search-patterns}

In many cases, when you're searching online, you may not know the
specific brand or company name for the product you need. Instead, you
have a general idea of the type of product you're looking for. For
instance, if you're in need of a pair of running shoes, you might simply
search for ``running shoes''.

\begin{figure}[!h]
  \centering
  \includegraphics[page=11, trim = 1.5cm 1cm 0.5cm 3cm, clip, width=\textwidth]{images/03 - Web_Information_Systems.pdf}
\end{figure}

When you enter a search term into Google, it aims to provide you with a
wide range of relevant results. Google takes pride in displaying a large
number of potential websites and useful links that address your search
query. For example, in the figure provided, the search term was
``fashion e-commerce'', and Google informs us that there are nearly 150
million possible sites that could meet this need.

\subsubsection{Importance of Top
  Rankings}\label{importance-of-top-rankings}

When users perform a search, they typically use common and direct terms.
As a company, it is crucial that your website's landing page appears in
the search results and that users click on the provided URL to visit
your site. However, it's important to understand that Google displays
search results in a text format, and users often need to click on the
links to determine if they are interested in the content. Simply reading
the brief description provided by Google is often not enough.

So, how patient are users when it comes to clicking through search
results? Studies have shown that if your company's website is ranked
below the 10th position on Google, only 10\% of users will click on the
link. This percentage decreases even further for results ranked between
10 and 20. In fact, only one in 100 users will click on these results.
This highlights the importance of achieving a high ranking in the search
results.

It's crucial to note that very few users venture beyond the first page
of search results. While it is possible to navigate to the second page,
the majority of users do not go that far. Therefore, it is essential to
aim for a top ranking to maximize online visibility and attract
potential clients.

\subsubsection{User Search Habits and
  Implications}\label{user-search-habits-and-implications}

It has been shown that Google is generally effective in providing search
results. However, for niche searches, the best fit result is often found
within the top 100 results rather than the top 10. This means that it is
important to be patient and take the time to go through the first 100
results, clicking on various landing pages to find the best one. This
level of thoroughness is typically expected from someone conducting
accurate market analysis. However, for the average user, especially a
potential retail customer, the priority is to find what they are looking
for quickly. If a result is not found within Google's top 10, there is a
tendency to assume that it does not exist. This is the conclusion that
most people draw when they cannot find a result within Google's top 10.

\subsubsection{Google's Role in Visibility and
  Monetization}\label{googles-role-in-visibility-and-monetization}

It is a common belief among retail customers that if a search result is
not within the top 10, it may not exist or be relevant. This behavior
can be conscious or unconscious, leading customers to give up on their
search after the first page of results. As a result, companies strive to
appear in Google's top 10 search results, as this is where most
visibility and traffic come from. Google capitalizes on this demand for
visibility by selling advertising space through its Google Ads platform.
Companies can pay to have their websites appear in the top results for
specific keywords.

However, this does not mean that the web provides perfect market
conditions for companies. In reality, companies must navigate through a
funnel controlled by Google, the dominant player in the search engine
market. Google profits from other companies' desire to rank high in its
search results. Whether this arrangement is fair or not is not for us to
judge in this course. Nevertheless, this system has been in place and
functioning for the past 20 years.

\subsection{Digital Divide and Marketing
  Strategies}\label{digital-divide-and-marketing-strategies}

\subsubsection{Challenges for Small
  Companies}\label{challenges-for-small-companies}

The digital divide poses a significant challenge for small companies,
especially when it comes to marketing. Limited financial resources make
it difficult for these companies to appear in search results. However,
even with a small budget, there are still ways to optimize Google
campaigns and gain visibility in a niche market. This niche can be
geographical or specific to less mainstream keywords, such as niche
issues or fashion e-commerce.

\subsubsection{Optimizing Google Campaigns and Niche
  Visibility}\label{optimizing-google-campaigns-and-niche-visibility}

In the realm of digital marketing, optimizing Google campaigns and
targeting niche audiences can help overcome the challenges of the
digital divide. While the concept may seem less mainstream, it is still
a viable strategy for businesses operating within a limited budget. By
focusing on a specific niche, companies can achieve visibility and make
an impact, even with limited resources.

\subsubsection{Marketplaces as a
  Solution}\label{marketplaces-as-a-solution}

\begin{figure}[!h]
  \centering
  \includegraphics[page=12, trim = 1.5cm 0.5cm 3cm 2.5cm, clip, width=\textwidth]{images/03 - Web_Information_Systems.pdf}
\end{figure}


Companies have responded to the opportunities and limitations of the
digital landscape in various ways, particularly when it comes to budget
constraints. One effective strategy they have employed is the creation
of marketplaces. Instead of solely focusing on establishing their own
online presence, companies choose to participate in existing
marketplaces. These marketplaces can gain significant popularity, such
as in the fashion industry where they may be associated with well-known
brands like Franciacorta or Ticino. The main concept behind
marketplaces is that multiple companies pool their resources to cover
the costs of online visibility. This approach offers clear advantages,
as companies can share the financial burden of establishing a strong
online presence, which can be quite substantial.

\paragraph{The Case of Yoox and Fashion
  Brands}\label{the-case-of-yux-and-fashion-brands}

E-commerce has brought both advantages and disadvantages to fashion brands. Take the case of Yoox, an e-commerce platform that sells products from established fashion brands. Initially, these brands overlooked the potential of online sales, believing that fashion products were not suitable for e-commerce. However, over time, consumer habits have changed, especially among younger demographics. It has become common for people to visit physical stores to try on products and then make their purchases online, seeking better deals or unique offerings.

While this shift in consumer behavior presents opportunities for fashion brands, there are also risks involved in not embracing technology and the online market. Many fashion companies now aspire to establish their own brand presence instead of relying on platforms like Yoox, as these platforms take a percentage of sales. By building their own online presence, brands can maintain greater control over their image and profit margins.

\subsection{Online Auctions and Dynamic
  Pricing}\label{online-auctions-and-dynamic-pricing}

\subsubsection{Concept of Online
  Auctions}\label{concept-of-online-auctions}

There are various ways to share the costs of visibility as a small
provider or seller. One effective method is to utilize online auction
services such as CDE. These platforms offer a range of online auction
options, some of which cater to specific product niches. These auctions
operate based on the concept of dynamic pricing, also known as options.

\subsubsection{Types of Auctions}\label{types-of-auctions}

\begin{figure}[!h]
  \centering
  \includegraphics[page=13, trim = 1.5cm 0cm 2cm 3cm, clip, width=\textwidth]{images/03 - Web_Information_Systems.pdf}
\end{figure}

\paragraph{Ascending Auctions}\label{ascending-auctions}

In the realm of dynamic pricing, there is a type known as ascending
auctions. These auctions are particularly suitable for goods or services
that are scarce. For example, if you have a piece of furniture with
historical value, it can be sold online through an ascending auction.
Since there are not many similar items available on the market, the
price for such a unique piece can vary greatly depending on its
condition.

In ascending auctions, the seller sets a base price and a minimum price.
Potential buyers then make offers, and the product is ultimately sold to
the highest bidder. It's important to act quickly in these auctions, as
there is a time limit for making offers. If you don't make a timely bid,
you risk losing the opportunity to purchase the item.

\paragraph{Descending or Dutch
  Auctions}\label{descending-or-dutch-auctions}

The rush to make offers in dynamic pricing is a key aspect of the
mechanism, especially when dealing with limited goods or services like
cars. In perfect market conditions, there would be an abundant and
unlimited supply of the good. However, in reality, this is not always
the case. To ensure you get the best price in such conditions, an
environment is created where people are inclined to make instinctive,
rather than entirely rational, offers. This leads to the price being
raised accordingly.

One type of auction that employs this strategy is the descending or
Dutch auction. In this type of auction, the vendor sets a very high
maximum price, which is then gradually decreased at regular time
intervals. The item is sold to the first client who stops the offer and
is willing to buy at the current price.

The sense of urgency remains, but the mechanism changes. The choice
between ascending and descending options depends on whether you, as a
vendor, are unwilling to sell below a certain base price or if the
market is unlikely to purchase above a specific amount.

\paragraph{Vickrey Auctions}\label{victory-auctions}

Another option is the Vickrey auction, where customers submit undisclosed offers within
a set timeframe, and the product is sold to the second highest bidder.
This type of auction is commonly used for tenders in the public
administration and by companies.

When companies need to purchase goods or services from a supplier, they
often use a document called a request for proposal (RFP). This document
is published and a selected number of suppliers are invited to
participate in an auction. The suppliers are required to submit another
document that explains how they will meet the requirements outlined in
the RFP. Additionally, they must include a separate envelope with a
price for the goods or services they are offering.

The reason why the product is sold to the second highest offer is to
prevent dumping. Dumping occurs when larger companies enter the market
by offering extremely low prices to gain market share. This can create a
monopoly and limit competition. By selling to the second highest offer,
companies can avoid dumping and promote a diverse market.

In some cases, auctions may prioritize the quality of the document that
explains the product or service being offered by the supplier. This
emphasizes the technical quality and ensures that the client receives
the best possible solution.


\subsection{E-Commerce Functionalities and Customer
  Profiling}\label{e-commerce-functionalities-and-customer-profiling}

\subsubsection{Advanced Functionalities}\label{advanced-functionalities}

\begin{figure}[!h]
  \centering
  \includegraphics[page=14, trim = 1.5cm 0cm 1.5cm 3cm, clip, width=\textwidth]{images/03 - Web_Information_Systems.pdf}
\end{figure}

E-commerce offers a range of advanced functionalities to meet the needs
of customers, whether it's through a marketplace or the commerce side of
a company. These functionalities are designed to enhance the customer
experience and provide efficient and effective solutions.

In e-commerce, there are several advanced functionalities that enhance
the customer experience. These include product configuration, pricing,
and the ability to place orders online. Customers have the convenience
of ordering and paying for their purchases on the website using various
payment methods such as credit, debit, or token-based payments. They can
also track the status of their orders online and access post-sale
services if they encounter any issues.

\subsubsection{Recommendation Systems and
  Personalization}\label{recommendation-systems-and-personalization}

\begin{figure}[!h]
  \centering
  \includegraphics[page=15, trim = 1cm 0cm 3cm 3cm, clip, width=\textwidth]{images/03 - Web_Information_Systems.pdf}
\end{figure}

As customers shop online, companies collect information about their
purchases to create customer profiles. This process, known as customer
profiling, helps companies understand their customers' preferences. For
example, if a customer clicks on a pastel-colored item, the
recommendation system will show them related products, such as sauces
for pasta. The recommendation system typically uses a combination of
collaborative filtering and content-based filtering, depending on the
company's business objectives.

When it comes to recommendation systems in e-commerce, personalization
is crucial, especially for companies that cater to different market
segments, such as the mid-market, low-end, and up-market. However, if a
company serving the up-market purchases an off-the-shelf recommendation
system that operates like Amazon's, there can be some challenges.

Typically, these off-the-shelf recommendation systems tend to suggest
mass-market products. This poses a risk for companies serving the
up-market because they may end up consistently recommending mid-market
products to customers who prefer up-market products. As a result,
instead of upselling and offering higher-priced products to increase
revenue, the company may unintentionally downsell and offer lower-priced
products.

In the best-case scenario, customers who are looking for up-market
products may simply ignore these recommendations because they are not
aligned with their preferences. However, in the worst-case scenario,
customers may find some interest in the mid-market recommendations and
choose those instead of the higher-priced up-market products. This can
lead to lower revenues and margins for the company.

\subsubsection{Growth Trends and Market
  Penetration}\label{growth-trends-and-market-penetration}

\begin{figure}[!h]
  \centering
  \includegraphics[page=16, trim = 2cm 0cm 3cm 4cm, clip, width=\textwidth]{images/03 - Web_Information_Systems.pdf}
\end{figure}

The extent to which e-commerce technologies are adopted depends on the
industry. In the case of e-fashion, it currently holds a market
penetration of approximately 20\%. This means that 80\% of sales still
take place in physical stores or through other channels. This is an
important observation because it highlights that e-commerce is just one
channel among many. Traditional companies must utilize multiple channels
if they wish to maintain their current revenue levels.

\subsection{Innovative E-Commerce
  Technologies}\label{innovative-e-commerce-technologies}

\subsubsection{Product Recommendations}\label{product-recommendations}

\begin{figure}[!h]
  \centering
  \includegraphics[page=17, trim = 0.5cm 0cm 0.5cm 1.5cm, clip, width=\textwidth]{images/03 - Web_Information_Systems.pdf}
\end{figure}

\begin{figure}[!h]
  \centering
  \includegraphics[page=18, width=\textwidth]{images/03 - Web_Information_Systems.pdf}
\end{figure}

To illustrate this point, let's consider advanced product recommendation
systems that leverage visual similarity. These systems are particularly
relevant in the fashion industry.

\subsubsection{Chatbots and Virtual
  Assistants}\label{chatbots-and-virtual-assistants}

\begin{figure}[!h]
  \centering
  \includegraphics[page=19, trim = 0.5cm 0cm 0.5cm 1.5cm, clip, width=\textwidth]{images/03 - Web_Information_Systems.pdf}
\end{figure}

Another way to enhance personalization in e-commerce is by utilizing
innovative technologies such as chatbots and virtual assistants. These
tools can serve as shopping assistants, providing personalized
recommendations based on the products you are interested in. For
example, when you are browsing a specific product, the chatbot or
virtual assistant can show you other similar products with different
customization options, allowing you to personalize your purchase. This
visual representation can help you make informed decisions and find the
perfect product that suits your preferences. Additionally, chatbots and
virtual assistants can provide real-time assistance, answering any
questions you may have and guiding you through the shopping process.
These technologies not only enhance the customer experience but also
streamline the purchasing journey, making it more efficient and
enjoyable.

To excel in the e-commerce industry, it is crucial to have a deep
understanding of customer preferences and be able to recommend suitable
products. This is where the application of AI can be highly
beneficial. By working tirelessly and becoming skilled at interpreting
and learning a customer's style, you can effectively propose products
that align with their preferences. This personalized approach can
greatly enhance the customer experience and increase satisfaction.

\subsubsection{Visual Search and Style
  Building}\label{visual-search-and-style-building}

\begin{figure}[!h]
  \centering
  \includegraphics[page=20, trim = 0.5cm 0.5cm 0.5cm 1cm, clip, width=\textwidth]{images/03 - Web_Information_Systems.pdf}
\end{figure}

Another innovative e-commerce technology that is worth exploring is
visual search. With visual search, customers can simply take a picture
of an item, such as a dress, and the recommendation system will search
for similar dresses and suggest them. This feature allows customers to
easily find products that match their desired style, making the shopping
experience more convenient and efficient. By incorporating visual search
into your e-commerce platform, you can provide a seamless and enjoyable
shopping experience for your customers.

\subsection{Conclusion and Transition to Part
  Three}\label{conclusion-and-transition-to-part-three}

Lastly, it is important to mention that your recommendation system can
also incorporate AI to personalize the style of
recommendations based on a specific client's purchases. This concludes
part two of the lecture, and in part three, we will delve into web
information systems, with a specific focus on supply chain management.

\section{Web Information Systems - III Lecture}

\subsection{Introduction to Supply Chain
  Management}\label{introduction-to-supply-chain-management}

Welcome to part three of our discussion on web information systems,
where we will be focusing on supply chain management. In the previous
sections, we explored the impact of e-commerce on retail customers and
how web information systems have revolutionized the way businesses
interact with their customers. However, it's important to note that the
web is not only used for e-commerce but also for e-business,
specifically in the context of coordinating exchanges and transactions
along the supply chain.

Supply chain management is a crucial aspect of e-business, where the web
serves as a platform for businesses to collaborate and execute
transactions with their partners along the supply chain. This innovative
application of web information systems has transformed the way
businesses operate and has led to significant improvements in efficiency
and effectiveness throughout the supply chain.

In the context of supply chain management in retail, the focus is on the
product that is ready to be sold to the end user. For example, when a
customer searches for ``rammy shoes'' and clicks on an e-fashion shop
that sells those shoes, they become the end user of the product.
However, in the realm of supply chain management, the interactions and
exchanges do not directly involve the end user.

\subsection{Understanding Supply Chains}\label{understanding-supply-chains}

\subsubsection{Definition and Components}\label{definition-and-components}

\begin{figure}[!h]
  \centering
  \includegraphics[page=22, trim = 2cm 7cm 2cm 1cm, clip, width=\textwidth]{images/03 - Web_Information_Systems.pdf}
\end{figure}

Supply chains consist of multiple businesses that work together to
produce and distribute a specific product. It's not just one company
involved, but rather a chain of companies that each handle different
aspects of production and distribution. These companies collaborate to
ensure that the product moves from raw materials to the final product
and reaches the end customer. In essence, a supply chain is a network of
companies that produce and sell the same product, with a shared focus on
serving the end users of that product.

\subsubsection{Complexity and Management}\label{complexity-and-management}

\begin{figure}[!h]
  \centering
  \includegraphics[page=23, trim = 1.5cm 3cm 4cm 1cm, clip, width=\textwidth]{images/03 - Web_Information_Systems.pdf}
\end{figure}

In supply chain management, companies closer to the end customer are
referred to as downstream companies, while those closer to raw materials
are known as upstream companies. The management of a supply chain is
facilitated by web information systems and tailored to the specific
needs of the product involved. Generally, a supply chain consists of a
cascade of companies operating in the same industry, from upstream to
downstream. These companies include suppliers, manufacturers,
distributors, retailers, and customers.

However, supply chains can be much more complex than this basic
structure. They can involve multiple levels of suppliers and customers
before reaching the final customers. Effective supply chain management
involves coordinating with suppliers and customers to achieve a common
goal: delivering the right product, at the right price, through the
right intermediary, in the right quantity, to the right customers, at
the right time.

\subsubsection{Financial Impact and
  Positioning}\label{financial-impact-and-positioning}

Supply chain management plays a crucial role in driving revenues and
profits for companies. The amount of money that companies can make along
the supply chain is directly influenced by effective supply chain
management. However, the financial impact and positioning of a company
within the supply chain can vary depending on its position and the
nature of the product.

While it is generally believed that companies closer to the customer,
downstream in the supply chain, tend to make more money, this is not
always the case. The profitability of a company depends on its strength
within the supply chain. This strength is determined by factors such as
bargaining power and the level of competition at that particular point
in the supply chain.

\subsubsection{Realistic View of Supply
  Chains}\label{realistic-view-of-supply-chains}

\begin{figure}[!h]
  \centering
  \includegraphics[page=24, trim = 3cm 0.5cm 3cm 4cm, clip, width=\textwidth]{images/03 - Web_Information_Systems.pdf}
\end{figure}

A realistic view of a supply chain involves multiple companies operating
along the same chain. In this diagram, the blue squares represent these
companies. For example, manufacturers within the supply chain may
compete against each other while also sharing suppliers and customers.
They compete to attract the attention of customers and suppliers,
ensuring they receive products at the right time to serve their
customers effectively.

\subsection{Supply Chain Management
  Process}\label{supply-chain-management-process}

\subsubsection{Continuous Learning Cycle}\label{continuous-learning-cycle}

\begin{figure}[!h]
  \centering
  \includegraphics[page=25, trim = 6cm 2cm 6cm 5cm, clip, width=\textwidth]{images/03 - Web_Information_Systems.pdf}
\end{figure}

Managing the supply chain is an ongoing process that requires continuous
learning. Companies must constantly strive to improve their ability to
manage the supply chain, as the world is constantly changing. This
continuous learning cycle allows companies to adapt and stay ahead. The
learning cycle is a repetitive process that ends at a certain point, but
then starts again when something changes and new knowledge needs to be
acquired.

\paragraph{Spend Analysis and
  Strategy}\label{spend-analysis-and-strategy}

To gain a better understanding of their supply chain, companies need to
go through a learning process when introducing a new product. The first
step in this process is spend analysis and strategy. This involves
examining how money is being spent and determining if it is being
allocated in the most effective way. Typically, the Pareto law is
applied at this stage. Companies evaluate their suppliers and the
materials they purchase, ranking them based on the amount of money
spent. Ideally, this ranking follows the Pareto principle, where 20\% of
the materials or goods account for 80\% of the spending. By focusing on
optimizing this 20\% of purchases, companies can achieve 80\% of the
benefits, as it covers the majority of their spending. This is the
essence of spend analysis.

\paragraph{Supplier Assessment and
  Sourcing}\label{supplier-assessment-and-sourcing}

In the supply chain management process, the first step is to assess the
supply market and determine if the company is purchasing from the best
suppliers. This involves creating a list of suppliers and evaluating
their position in the market. Strategic planning is then done to
determine if any suppliers should be replaced with better options. The
company may also consider experimenting with new suppliers who are
market leaders.

Once the best suppliers have been identified, the next step is sourcing.
This involves negotiating the total cost with potential suppliers. The
company approaches the supplier and discusses the amount they currently
spend on the type of goods or materials provided by the supplier.

\paragraph{Requirements Management}\label{requirements-management}

In our supply chain management process, one of the key steps is spend
analysis and strategy. Each year, we analyze our spending and develop a
strategy based on that information. As a supplier, if you can offer us a
discount or integrate our information systems, it would make it more
convenient for us to work with you. For example, if we have access to
your stock levels, we can place orders with confidence, knowing that you
can meet our requirements in a timely manner. We also prioritize
electronic transactions to ensure accuracy and efficiency.

In some cases, we may even reach a point where we have automated
processes, such as automatic order placement through MRP (Material
Requirements Planning). Of course, this would only be possible if the
supplier has the capacity to fulfill the orders.

Another important aspect of our relationship with suppliers is managing
new requirements. The market is constantly changing, and we need to
adapt to the evolving needs of our end customers. We work closely with
our suppliers to address these requirements together. We value
flexibility and the ability to respond to emerging needs, rather than
rigid standardization. This collaborative approach allows both parties
to continuously grow and stay competitive in the market.

In fact, our relationship with suppliers can be so strong that we even
collaborate on designing new products or services together. This level
of partnership and cooperation is what we aim for in our requirements
management process.

\paragraph{Supplier Relationship
  Management}\label{supplier-relationship-management}

In order to effectively manage supplier relationships, it is important
to focus on step number three: requirements management. This step should
be reserved for suppliers who have proven their loyalty and ability to
support your business. By concentrating your purchases with these
dependable suppliers, you can establish a strong working relationship
over time.

However, it is crucial to recognize that suppliers may close, be sold,
or acquired by other companies, regardless of their loyalty. In such
cases, it becomes necessary to find new suppliers who can meet your
requirements. To ensure a smooth transition, it is essential to engage
in supplier relationship management.

Supplier relationship management is similar to a CRM package, but
specifically designed for suppliers. It provides a platform where
suppliers can access a workspace, download contracts and RFPs,
participate in tenders, and receive notifications about upcoming
opportunities. They can also manage their contracts, payments, and
undergo continuous evaluation and qualification processes.

By actively managing supplier relationships through this system,
suppliers have the opportunity to grow and improve their performance.
Positive evaluations can give them a competitive edge in future tenders,
ensuring a mutually beneficial partnership.

\subsection{Technology in Supply Chain
  Management}\label{technology-in-supply-chain-management}

In addition to attending courses online, companies can also provide
training programs for their suppliers. Through e-learning platforms,
suppliers can access courses and gain knowledge about the company's
needs before being approached for new projects. This allows suppliers to
grow and improve their capabilities, making them more suitable for
collaboration with the company. Certification and qualification
processes further ensure that suppliers meet the company's standards.

All of these processes are managed through technology, which includes
analytics, spend analysis, and strategy. The negotiation phase involves
sourcing, and requirements management is done collaboratively.
Transaction support is provided through e-commerce, and as the
relationship with suppliers develops, the platform evolves to include
CRM-like features for managing supplier relationships. This
comprehensive technology platform enables efficient and effective
collaboration between companies and their suppliers.


\end{document}